\section{A rational functional and a positive measure}\label{sec:functional}
\source{Construction: Fauzan \cite[(2.2)--(2.3)]{fauzan} and Lai
\cite[Definition 4.2]{lai}. The presentation follows the functional
viewpoint in Persiflage \cite{persiflage}.}
\subsection{Defining the functional carefully}
The notation involves three different variables. The letter $t$ is the
argument of a rational function. The letter $X$ is a parameter that will
eventually be set equal to $\xi$. The letter $y$ will be an integration
variable. Keeping these roles separate makes the later matrix calculation
much easier to follow.

Let $\mathcal R$ be the vector space over $\Q$ of rational functions in
$\Q(t)$ whose finite poles, if any, are simple and belong to
$\{-1^2,-2^2,-3^2,\ldots\}$. Every $R\in\mathcal R$ has a unique expression
\begin{equation}\label{eq:partial}
 R(t)=\sum_{e=0}^{d}p_et^e+\sum_{j\in J}\frac{c_j}{t+j^2},
 \qquad p_e,c_j\in\Q,
\end{equation}
where $J$ is finite. Polynomial division gives the polynomial part.
At a pole, its coefficient is
$c_j=\left.(t+j^2)R(t)\right|_{t=-j^2}$.
Subtracting these simple fractions leaves a rational function with no finite
poles, hence a polynomial. These observations also prove uniqueness
(with zero terms omitted).

Define the Bernoulli numbers by
\[
 \frac{u}{e^u-1}=\sum_{k\geq0}B_k\frac{u^k}{k!}.
\]
Multiplication by the exponential series determines the $B_k$ recursively
as rational numbers. In particular $B_0=1$, $B_1=-1/2$, $B_2=1/6$,
and $B_4=-1/30$. Also put
\[
 \Hfive{j}=\sum_{v=1}^j v^{-5},\qquad \Hfive{0}=0.
\]
Prescribe a $\Q$-linear functional $\mu_X:\mathcal R\to\Q[X]$ by
\begin{align}
 \mu_X(t^e)&=
 \frac{(-1)^eB_{2e+2}(2e+3)(2e+4)(2e+5)}{24},
 &&e\in\Z_{\geq0},\label{eq:mu-moment}\\
 \mu_X\!\left(\frac1{t+j^2}\right)&=
 j^4\bigl(X-\Hfive{j}\bigr)-\frac14+\frac1{2j},
 &&j\in\Z_{>0}.\label{eq:mu-pole}
\end{align}
For a general $R$, apply these formulas to its decomposition and add with
the coefficients $p_e,c_j$. This is the meaning of extending by linearity.
It gives a well-defined affine polynomial in $X$: only the simple
fractions contribute a term involving $X$.

For example,
\[
 \mu_X(1)=\frac5{12},\qquad
 \mu_X\!\left(\frac1{t+1}\right)=X-\frac34.
\]
Since $t/(t+1)=1-1/(t+1)$, linearity gives
$\mu_X(t/(t+1))=7/6-X$. There is no rule allowing us to multiply
functional values; $\mu_X$ is linear, not multiplicative.

Finally, $\mu_\xi(R)$ means evaluate the polynomial $\mu_X(R)$ at $X=\xi$.
It does not mean substitute $\xi$ for $t$ in the rational function.

\subsection{Why the polynomial values are moments}
\source{Pr\'evost \cite[Theorems 1--2]{prevost};
Fauzan \cite[Proposition 2.2]{fauzan}; Lai \cite[Lemma 4.5]{lai}.}
Define, for $y>0$,
\begin{equation}\label{eq:weight}
 w(y)=\frac{(2\pi)^4y^5}{12}
       \sum_{\ell\geq1}\ell^4e^{-2\pi\ell y}.
\end{equation}
Every summand is positive. The series converges for each $y>0$, so
$w(y)>0$. We will prove
\begin{equation}\label{eq:representation}
 \boxed{\mu_\xi(R)=\int_0^\infty R(y^2)w(y)\dd y
        \quad(R\in\mathcal R).}
\end{equation}
This is the bridge between rational algebra and real integration.
First we check that the integrals make sense.

Put $f(y)=(e^{2\pi y}-1)^{-1}$. Its geometric series may be differentiated
term by term on $y\geq\varepsilon>0$, where the derivative series converge
uniformly. Four differentiations show that
\[
 w(y)=\frac{y^5}{12}f^{(4)}(y).
\]
Near zero, $f(y)=1/(2\pi y)+g(y)$ with $g$ analytic at zero. Consequently
$f^{(4)}(y)=24/(2\pi y^5)+O(1)$ and $w(y)\to1/\pi$.
For $y\geq1$, factor $e^{-2\pi y}$ out of the series in \eqref{eq:weight};
the remaining sum is bounded by
$\sum_{\ell\geq1}\ell^4e^{-2\pi(\ell-1)}<\infty$.
Thus $w(y)=O(y^5e^{-2\pi y})$ at infinity.

An admissible $R(y^2)$ is bounded near zero, because its poles lie away from
$[0,\infty)$. It grows at most polynomially as $y\to\infty$.
These two bounds on $w$ prove absolute convergence in
\eqref{eq:representation}.

For $R(t)=t^e$, the summands of the integral are nonnegative. Tonelli's
theorem allows us to interchange the sum and integral:
\begin{align*}
 \int_0^\infty y^{2e}w(y)\dd y
 &=\frac{(2\pi)^4}{12}\sum_{\ell\geq1}\ell^4
      \int_0^\infty y^{2e+5}e^{-2\pi\ell y}\dd y\\
 &=\frac{(2e+5)!}{12(2\pi)^{2e+2}}
      \sum_{\ell\geq1}\frac1{\ell^{2e+2}}.
\end{align*}
Here $\int_0^\infty y^m e^{-by}\dd y=m!/b^{m+1}$ for $b>0$:
substitution reduces this to $b=1$, and repeated integration by parts gives
$m!$. Euler's formula, proved in Appendix A, is
\[
 \zeta(2e+2)=
 \frac{(-1)^eB_{2e+2}(2\pi)^{2e+2}}{2(2e+2)!}.
\]
Substituting it above yields
\[
 \int_0^\infty y^{2e}w(y)\dd y
 =\frac{(-1)^eB_{2e+2}}{24}\,
       \frac{(2e+5)!}{(2e+2)!}
 =\mu_X(t^e).
\]
The factorial quotient is exactly $(2e+3)(2e+4)(2e+5)$.
The equality is independent of $X$, as the definition requires.

\subsection{Why the pole values are the same integrals}
\source{The Stieltjes representation is Pr\'evost's \cite[Theorem 1]{prevost}.
The fifth-order calculation is also in Fauzan \cite[Proposition 2.2]{fauzan}
and Lai \cite[Lemma 4.6]{lai}.}
The simple fractions need a different calculation. We do it for a real
number $a>0$, then specialize to a positive integer. Define
\[
 q_a(y)=\frac{y^5}{y^2+a^2}
       =y^3-a^2y+a^4\frac{y}{y^2+a^2}.
\]
The polynomial terms disappear after four derivatives. Since
$y/(y^2+a^2)=\operatorname{Im}(a-iy)^{-1}$, and
$\frac{d^4}{dy^4}(a-iy)^{-1}=24(a-iy)^{-5}$, we obtain
\[
 q_a^{(4)}(y)=24a^4\operatorname{Im}(a-iy)^{-5}.
\]
Using $w=y^5f^{(4)}/12$, four integrations by parts give
\begin{align*}
 \int_0^\infty\frac{w(y)}{y^2+a^2}\dd y
 &=\frac1{12}\int_0^\infty q_a(y)f^{(4)}(y)\dd y\\
 &=\frac1{12}\int_0^\infty q_a^{(4)}(y)f(y)\dd y\\
 &=2a^4\int_0^\infty
     \frac{\operatorname{Im}(a-iy)^{-5}}{e^{2\pi y}-1}\dd y.
\end{align*}
There are no missing boundary terms. First work on $[\varepsilon,L]$;
the boundary expression is
$q_af'''-q_a'f''+q_a''f'-q_a'''f$.
For $0\leq r\leq3$, at zero we have
\[
 q_a^{(r)}(y)=O(y^{5-r}),\qquad
 f^{(3-r)}(y)=O(y^{r-4}),
\]
so every boundary product is $O(y)$ and tends to zero. At infinity the
derivatives of $f$ decay exponentially and the other factors have at most
polynomial growth. Also $q_a^{(4)}(y)=O(y)$ at zero, so
$q_a^{(4)}(y)f(y)=O(1)$ there. The final integral converges at both ends,
and we may let $\varepsilon\downarrow0$ and $L\to\infty$.

For $a>0$, define $\zeta(5,a)=\sum_{v=0}^\infty(a+v)^{-5}$.
The particular Hurwitz identity needed here is
\begin{equation}\label{eq:hermite}
 \zeta(5,a)=\frac1{2a^5}+\frac1{4a^4}
 +2\int_0^\infty
     \frac{\operatorname{Im}(a-iy)^{-5}}{e^{2\pi y}-1}\dd y.
\end{equation}
Appendix A derives this directly from an elementary kernel identity.
It gives
\[
 \int_0^\infty\frac{w(y)}{y^2+a^2}\dd y
 =a^4\zeta(5,a)-\frac1{2a}-\frac14.
\]
For an integer $j\geq1$, the series gives
$\zeta(5,j)=\xi-\Hfive{j}+j^{-5}$. Therefore
\[
 \int_0^\infty\frac{w(y)}{y^2+j^2}\dd y
 =j^4(\xi-\Hfive{j})-\frac14+\frac1{2j}
 =\mu_\xi\!\left(\frac1{t+j^2}\right).
\]
We have checked both types of term in \eqref{eq:partial}.
Applying the formulas to its finitely many terms and using linearity
of integration proves \eqref{eq:representation} for every $R\in\mathcal R$.

\subsection{One transform, two kinds of data}
\source{This perspective is guided by Persiflage \cite{persiflage}.
The underlying transform and orthogonal-polynomial identities are
Pr\'evost's \cite[Section 2 and Theorems 1--2]{prevost}.}

It is helpful to see where the zeta value enters.
Change variables \(t=y^2\) and write
\[
 d\nu(t)=\frac{w(\sqrt t)}{2\sqrt t}\,dt,\qquad
 \Phi(z)=\int_0^\infty\frac{d\nu(t)}{z+t}\quad(z>0).
\]
The moments \(m_j=\int t^j\,d\nu(t)=\mu_X(t^j)\) are rational.
On the other hand, the pole formula gives
\[
 \Phi(a^2)=a^4\bigl(\xi-H_a^{(5)}\bigr)-\frac14+\frac1{2a}
 \qquad(a\in\Z_{>0}).
\]
Thus rational approximation to a transform value gives rational
approximation to \(\xi\).

The moments also describe \(\Phi\) for large \(z\).
The finite geometric identity gives
\[
 \Phi(z)=\sum_{j=0}^{m-1}\frac{(-1)^jm_j}{z^{j+1}}
  +\frac{(-1)^m}{z^m}\int_0^\infty\frac{t^m}{z+t}\,d\nu(t).
\]
The last term has absolute value at most \(m_m z^{-m-1}\).
This is an asymptotic expansion for each fixed \(m\); we have not asserted
convergence of the infinite moment series.

At \(X=\xi\), the two clauses defining \(\mu_X\) describe samples and
asymptotic coefficients of one transform. This is the bridge to
Pr\'evost's Pad\'e approximation, emphasized by Persiflage.
Our proof uses a rationally modified measure and its determinant;
no quantitative Pad\'e estimate is needed. The question is how much
real smallness survives when the determinant's coefficients are made
integral.

\section[Modifying the measure and its determinant]{Modifying the measure and taking its determinant}\label{sec:determinant}
\source{Fauzan \cite[Section 2.1]{fauzan}; Persiflage
\cite[Hankel Determinants]{persiflage}; Lai \cite[Section 5]{lai}.}
\subsection{The modifier and its affine moments}
For a positive integer $n$, choose
\[
 K=40n,\qquad N=3n,\qquad h=K-N=37n,
 \qquad D_m(t)=\prod_{j=1}^m(t+j^2).
\]
The parameters have been chosen to make the later estimates work.
For now, the features that matter are positivity and cancellation. Define
\[
 A_K(t)=\frac{D_N(t)^6}{D_K(t)}
       =\frac{D_N(t)^5}{T_K(t)},\qquad
 T_K(t)=\prod_{k=N+1}^K(t+k^2).
\]
The cancellation removes the first $N$ possible poles. The $h$ remaining
roots of $T_K$ are distinct negative integer squares, so multiplying
$A_K$ by any polynomial leaves a function in $\mathcal R$.
Moreover, $A_K(y^2)>0$ for every $y>0$.

It is useful to name the pairing before choosing its coordinates:
\begin{equation}\label{eq:bilinear-form}
 \mathcal B_{K,X}(p,q)=\mu_X(A_Kpq),
 \qquad p,q\in\Q[t]_{<h}.
\end{equation}
The monomials below and the later bases adapted to each prime are
coordinates for this one form. The real estimate uses its integral
representation. This common starting point is explicit in
Fauzan \cite[Sections 2.3--2.4]{fauzan}.

For $0\leq i,j<h$ put $R_{ij}(t)=A_K(t)t^{i+j}$ and define
\begin{equation}\label{eq:G}
 G_K(X)=\bigl[\mu_X(R_{ij})\bigr]_{0\leq i,j<h},
 \qquad \Delta_K(X)=\det G_K(X).
\end{equation}
The dependence on $i+j$ makes this a Hankel matrix: its entries are constant
along each anti-diagonal.
To see explicitly why these are affine entries, take partial fractions:
\[
 R_{ij}(t)=P_{ij}(t)+\sum_{k=N+1}^K\frac{c_k^{ij}}{t+k^2},
 \qquad
 c_k^{ij}=\frac{D_N(-k^2)^5(-k^2)^{i+j}}{T_K'(-k^2)}.
\]
The polynomial part may be zero; it must not otherwise be discarded.
Applying the functional gives
\[
 G_K(X)_{ij}=\mu_X(P_{ij})+
 \sum_{k=N+1}^K c_k^{ij}
 \left[k^4(X-\Hfive{k})-\frac14+\frac1{2k}\right].
\]
Thus every entry is $a_{ij}X+b_{ij}$ with rational $a_{ij},b_{ij}$.
Each term of the determinant expansion is a product of \(h\) affine
entries. Hence $\Delta_K\in\Q[X]$ and $\deg\Delta_K\leq37n$.
Section~\ref{sec:exact-degree} proves exact degree $h$, but the upper bound suffices.

\subsection{Evaluating the matrix at \texorpdfstring{\(\xi\)}{xi}}
\source{The positive Gram interpretation is explicit in Fauzan
\cite[Proposition 2.2]{fauzan}, Persiflage \cite{persiflage},
and Lai \cite[Proposition 5.6]{lai}.}
We can now pass from a symbolic matrix to a matrix of integrals one entry
at a time. With $R_{ij}$ as above,
\begin{align*}
 G_K(\xi)_{ij}
 &=\left.\mu_X(R_{ij})\right|_{X=\xi}\\
 &=\mu_\xi(R_{ij})\\
 &=\int_0^\infty R_{ij}(y^2)w(y)\dd y\\
 &=\int_0^\infty y^{2i}y^{2j}A_K(y^2)w(y)\dd y.
\end{align*}
The third equality is the representation proved in Section~\ref{sec:functional}.
The last is substitution: $(y^2)^{i+j}=y^{2i}y^{2j}$.
Absolute convergence was already checked for every admissible rational
function, including these particular $R_{ij}$.

For a compact description, set
\[
 \dd\nu_K(y)=A_K(y^2)w(y)\dd y,\qquad
 v(y)=(1,y^2,\ldots,y^{2h-2})^{\mathsf T}.
\]
Then $G_K(\xi)=\int_0^\infty v(y)v(y)^{\mathsf T}\dd\nu_K(y)$,
where the integral is taken \emph{entry by entry}. Its $(i,j)$ entry is
exactly the last integral above.

Let $c=(c_0,\ldots,c_{h-1})^{\mathsf T}\in\R^h$ and
$q(t)=\sum_i c_it^i$. Distributing the finite sums gives
\[
 c^{\mathsf T}G_K(\xi)c
 =\int_0^\infty\left(\sum_i c_iy^{2i}\right)^2\dd\nu_K(y)
 =\int_0^\infty q(y^2)^2\dd\nu_K(y).
\]
Only finite linearity is used here. Real coefficients $c_i$ cause no
difficulty with the originally $\Q$-linear definition of $\mu_X$.

If $c\ne0$, then $q$ is a nonzero polynomial. Choose $y_0>0$ that is not
a root of $q(y^2)$. By continuity, $q(y^2)^2$ is bounded below by a positive
number on a small interval about $y_0$. The density of $\nu_K$ is positive
and continuous there, so the integral is strictly positive.

Thus $G_K(\xi)$ is positive definite. To see why its determinant is positive,
orthogonalize $1,y^2,\ldots,y^{2h-2}$ by Gram--Schmidt without normalizing.
The resulting orthogonal functions are monic polynomials in $y^2$, so each
is nonzero and has positive squared norm. The change of basis is triangular
with diagonal entries one. The determinant is therefore the product of
these positive squared norms: $\Delta_K(\xi)>0$.

The matrix integral is not a rule for moving a determinant inside an
integral. Each $v(y)v(y)^{\mathsf T}$ has rank one. The multiple integral
for the determinant comes from Andr\'eief's identity, proved in the analytic section.

The integral and the affine dependence can now be displayed together.
For rational polynomials \(p,q\) of degree less than \(h\), put
\[
 d_k=\frac{k^4D_N(-k^2)^5}{T_K'(-k^2)}
 \qquad(N<k\le K).
\]
The partial fractions already computed give the exact identity
\begin{equation}\label{eq:gram-residue}
 \mathcal B_{K,X}(p,q)
 =\int_0^\infty p(y^2)q(y^2)\dd\nu_K(y)
 +(X-\xi)\sum_{k=N+1}^K d_kp(-k^2)q(-k^2).
\end{equation}
Indeed the polynomial part is independent of \(X\), while the change
in the value at the pole \(-k^2\) is \(k^4(X-\xi)\).
The residue sum is the slope; polynomial quotients still contribute to
the constant part. This explains the later arithmetic correction.
The identity combines Fauzan's Gram and residue formulas
\cite[Sections 2.3--2.4]{fauzan}.

\subsection{Putting both estimates on the same scale}
Fauzan uses the positive rational scalar
\begin{equation}\label{eq:normalization}
 S_K=\frac{(K!)^{2h}4^{h-1}}
 {(N!)^{12h}\prod_{i=1}^{h-1}((2i)!)^2},
 \qquad F_K(X)=S_K\Delta_K(X).
\end{equation}
This is his equation (2.5) \cite{fauzan}, also retained in
\texttt{Apery/Construction.lean} \cite{firsching}.
It follows at once that
\[
 F_K\in\Q[X],\qquad \deg F_K\leq h,\qquad F_K(\xi)>0.
\]
The rescaling is useful because factorial valuations admit exact formulas,
and factorial sizes admit Stirling estimates. It does not, on its own,
make every coefficient integral.

The goal itself is independent of rational scale. Let \(c(A)>0\)
be the rational content of a nonzero \(A\in\Q[X]\): dividing by \(c(A)\)
makes its coefficients integers with greatest common divisor one.
If \(v_p^{\mathrm G}(A)\) denotes the minimum coefficient valuation,
then \(c(A)=\prod_p p^{v_p^{\mathrm G}(A)}\), by taking a common
denominator and then the greatest common divisor. For \(A(\xi)\ne0\), set
\begin{equation}\label{eq:primitive-size}
 \mathcal J_\xi(A)
 =\log|A(\xi)|-\sum_p v_p^{\mathrm G}(A)\log p
 =\log\left|\frac{A(\xi)}{c(A)}\right|.
\end{equation}
For \(r\in\Q^\times\), the identity
\(\log|r|=\sum_p v_p(r)\log p\) proves
\(\mathcal J_\xi(rA)=\mathcal J_\xi(A)\).
The real change in size caused by a rational rescaling is exactly
balanced by its change at the finite primes.
Fauzan already uses primitive normalization
\cite[Appendix C.3]{fauzan}. We choose \(S_K\) for convenient estimates,
then construct an explicit multiplier from certified prime bounds.
Together those estimates bound \(\mathcal J_\xi\).
